\subsection{Machine Learning and AI Research}

Machine-learning and AI conclusions draw on statistical, computational, system, and social evidence; the study design must identify the authority supporting each conclusion.

\begin{process} \label{proc:ml-ai-study-protocol} \hlt{Machine-Learning and AI Study Protocol}\\
Specify an empirical study before interpreting its scores.
\begin{enumerate}[label=(\roman*)]
\setlength{\itemsep}{0pt}
\item \hlt{Question}: state the population, task, context of use, intervention or comparison, and claimed outcome.
\item \hlt{Claim type}: classify each target as theoretical, empirical, causal, computational, system, or normative.
\item \hlt{Data}: record provenance, version, license, collection, consent or access status, labels, exclusions, and preprocessing.
\item \hlt{Evaluation boundary}: protect test data from training, model selection, prompt design, preprocessing, and human iteration.
\item \hlt{Comparators}: choose relevant baselines, ablations, and resource-matched alternatives before inspecting final results.
\item \hlt{Protocol}: fix splits, metrics, hyperparameter selection, seeds, stopping rules, hardware, software, and compute budget.
\item \hlt{Uncertainty}: repeat stochastic runs and report how intervals, variation, or significance measures are calculated.
\item \hlt{Stress tests}: examine subgroups, shift, robustness, failure modes, sensitivity, and plausible misuse within scope.
\item \hlt{Artifacts}: preserve code, data handles, checkpoints, configurations, logs, and command-to-result mappings.
\item \hlt{Report}: separate observed results from interpretation and state limitations, costs, risks, and unsupported generalizations.
\end{enumerate}
The 2026 NeurIPS checklist and its reproducibility program emphasize assumptions, complete proofs, experimental detail, uncertainty, artifacts, ethics, and limitations \citep{neurips2026checklist,pineau2021reproducibility}.
\end{process}

\begin{process} \label{proc:ml-artifact-documentation} \hlt{Dataset and Model Documentation}\\
Publish documentation alongside every durable ML artifact.
\begin{enumerate}[label=(\roman*)]
\setlength{\itemsep}{0pt}
\item \hlt{Dataset}: record motivation, version, provenance, composition, collection, preprocessing, labeling, license, intended uses, restrictions, and maintenance.
\item \hlt{Model}: record identity, version, intended and excluded uses, training data, evaluation protocol, subgroup results, and limitations.
\item \hlt{System}: record interfaces, human roles, deployment assumptions, monitoring, update policy, safeguards, and rollback conditions.
\end{enumerate}
Dataset datasheets and model cards provide complementary documentation patterns for provenance, intended use, evaluation, and limitations \citep{gebru2021datasheets,mitchell2019modelcards}.
\end{process}

\begin{principle} \label{princ:contextual-ai-risk} \hlt{Contextual AI Risk Assessment}\\
Accuracy alone is not a complete system claim. For intended context, select and justify tests for the following applicable characteristics.
\begin{enumerate}[label=(\roman*)]
\setlength{\itemsep}{0pt}
\item validity and reliability;
\item safety;
\item security and resilience;
\item accountability and transparency;
\item explainability and interpretability;
\item privacy;
\item harmful-bias management.
\end{enumerate}
These characteristics can conflict and require documented tradeoffs. \\
his versioned vocabulary follows the voluntary NIST AI RMF 1.0 \citep{tabassi2023airmf}.
\end{principle}

\begin{warning} \label{warn:ml-evaluation-leakage} \hlt{Adaptive Evaluation and Leakage}\\
Any choice influenced by protected evaluation data makes the final score adaptive. Repeated leaderboard queries, test-set debugging, manual example selection, and unreported prompt tuning can leak evaluation information.\\
Reserve a final assessment boundary; otherwise disclose the adaptation and narrow the claim.
\end{warning}

\begin{warning} \label{warn:ml-benchmark-overclaiming} \hlt{Benchmark Overclaiming}\\
A higher score supports a claim only under the reported dataset, metric, resources, selection procedure, and uncertainty.\\
It does not establish practical usefulness, causal benefit, robustness, fairness, safety, or superiority in another operating environment without matching evidence.
\end{warning}
